\subsection{The smooth critical curve of transverse type \texorpdfstring{$A_3$}{A3}}

\begin{statement}{P4R.14}{Proved}
Let $S$ be an actual reduced support in Part~4.2. At a point
where the entire reduced critical locus of its local defining
function is one smooth curve, with generic transverse type
$A_3$, holomorphic coordinates put that function in one of
the forms
\[
 X^2+Y^4,\qquad X(Y^2+XZ).                              \tag{P4R.14.1}
\]
The normalization germs are respectively two smooth branches,
or a smooth plane and an $A_1$ quadric. No positive-genus
exceptional curve occurs over such a point.
\end{statement}
\begin{proof}
Choose a coordinate $z$ on the critical curve and normal
coordinates $x,y$. Its quadratic transverse Hessian has generic
rank one. Diagonalization of the symmetric matrix over
$\mathbb C\{z\}$, as in P4.28, gives the quadratic term
$h(z)x^2$. The cubic coefficient on the null direction
vanishes identically, by the generic transverse type. With
weights two for $x$ and one for $y$, the weight-four terms are
\[
 h(z)x^2+a(z)xy^2+b(z)y^4,\qquad
 \Delta(z)=4h(z)b(z)-a(z)^2.                             \tag{P4R.14.2}
\]
Every remaining term has larger weight and $\Delta$ is
nonzero generically.

We first prove that $\Delta$ is a unit. A zero persists,
counted with multiplicity, under small independent constant
perturbations of $h,a,b$, by the argument principle on a
small parameter disc. Generic such perturbations make the
zeros simple and disjoint from $h=0$. Indeed $h=\Delta=0$
also requires $a=0$, which the independent parameters avoid;
where $h\ne0$, the derivative of $\Delta$ in the $b$
parameter is $4h\ne0$. These perturbations preserve the
critical axis and its transverse $A_3$ germ near the fixed
boundary. At a simple zero with $h\ne0$, the parametrized
Morse lemma gives
\[
 X^2+y^4U(y,z),\qquad U(0,0)=0,\quad U_z(0,0)\ne0.
\]
The actual coordinate $Z=U(y,z)$ gives $X^2+y^4Z$.
P4.16 supplies a compact spherical class with nonzero square
in its smoothing. The boundary comparison of P4.28 applies
here: near the boundary critical curve, the parametrized
Morse lemma and a fourth root of a unit give $X^2+Y^4$;
the tangential boundary coordinate makes the nonzero levels
submersive. Proper boundary transport therefore carries the
sphere to the actual nearby torus. Its universal cover is
$\mathbb R^4$, so a sphere has zero homology class and zero
square, a contradiction. Thus $\Delta$ is a unit.

For $h(0)\ne0$, Morse splitting and a fourth root of the
residual unit give the first form in (P4R.14.1). Suppose
$h(0)=0$. Then $a(0)\ne0$. The holomorphic function
\[
 \frac{f(y^2v,y,z)}{y^4}
\]
restricts at $y=0$ to $h(z)v^2+a(z)v+b(z)$. Its finite
root $-b(0)/a(0)$ at the origin is simple. The holomorphic
implicit-function theorem gives a root $v=v(y,z)$, and
division by the monic equation of its graph gives
\[
 X=x-y^2v(y,z),\qquad f=Xg.
\]
On $X=0$, the function $g$ is $y^2$ times a unit. Its
parameter Morse centre in $y$ is therefore zero there and
its critical value is zero. The parametrized Morse lemma,
with parameters $X,z$, gives the exact form
\[
 f=X\bigl(Y^2+XH(X,z)\bigr).                            \tag{P4R.14.3}
\]
The remaining critical value is divisible by $X$ by its
vanishing on $X=0$. The function $H(0,z)$ is not identically
zero, since the generic transverse Hessian has rank one;
at the point under consideration $H(0,0)=0$.

Put $G=2H+XH_X$. Off $X=0$, the critical equations of
(P4R.14.3) are $Y=0,H_z=0,G=0$. We prove $H_z(0,0)\ne0$.
Otherwise $H_z$ is a nonzero plane germ not divisible by
$X$, because $H(0,z)$ has finite order at least two.
No component of $H_z=0$ is contained in $G=0$: it would
give an additional critical-curve branch off $X=0$, whereas
the entire critical locus is the stated smooth axis. Hence
$(H_z,G)$ is an isolated complete intersection of positive
intersection multiplicity.

Replace $H$ by $H-\eta z-\varepsilon$. Its plane critical
equations are
\[
 H_z-\eta=0,\qquad G-2\eta z-2\varepsilon=0.              \tag{P4R.14.4}
\]
Their solutions are the fibres of the holomorphic map
\[
 \Psi(X,z)=\left(H_z,\frac{G-2zH_z}{2}\right).
\]
Its zero fibre has the positive local degree just proved.
Choose a small compact representative with no boundary zero.
That degree is unchanged on a sufficiently small open
parameter ball. The universal equations in (P4R.14.4) have
parameter Jacobian determinant two, so regular parameter
values give a nonempty finite set of transverse solutions.
We also choose such a value outside the images under $\Psi$
of $X=0$ and $H_X=0$. These exceptional images have real
four-dimensional measure zero: each source curve has
countably many smooth real two-dimensional strata, and the
smooth image of each has measure zero. The germ $H_X$ is
not identically zero. Indeed in that case $H$ would depend
only on $z$ and both $H_z$ and $G=2H$ would vanish on the
whole line $z=0$, contrary to the isolated intersection.
The two exceptional images cannot exhaust the open ball
on which the positive degree is retained. This argument
uses no analyticity assertion for the image of a nonproper
map.

At each chosen solution, the critical value is
\[
 f_{\eta,\varepsilon}=-\frac{X^3H_X}{2}\ne0.              \tag{P4R.14.5}
\]
For the indicated order of the two equations, put
\[
 J=\det\frac{\partial(G-2\eta z-2\varepsilon,\ H_z-\eta)}
                  {\partial(X,z)}.
\]
The $Y$-Hessian entry is $2X$, and the remaining Hessian
block at a critical point is
\[
 \begin{pmatrix}
 X(3H_X+XH_{XX})&X^2H_{Xz}\\
 X^2H_{Xz}&X^2H_{zz}
 \end{pmatrix}.
\]
Its determinant is $X^3J$, and the full Hessian determinant
is $2X^4J\ne0$. Thus these are isolated holomorphic Morse
points with nonzero critical values. Near the fixed boundary,
$H(0,z)$ remains nonzero under the small perturbation, so the
critical axis retains its transverse $A_3$ type. The same
proper boundary comparison transports a Morse vanishing
sphere, of square $-2$, into an actual nearby torus. This is
impossible. We have proved $H_z(0,0)\ne0$.

The holomorphic coordinate $Z=H(X,z)$, keeping $X,Y$, now
gives the second form of (P4R.14.1). Its normalization has
the plane $X=0$ and the normal quadric $Y^2+XZ=0$. The
quadric's minimal resolution has one rational $(-2)$-curve.
The first form factors as $(X+iY^2)(X-iY^2)$ and its two
normalization branches are smooth. This proves the stated
normalization conclusions.

The distinction between a pole coefficient and this
torus-smoothing conclusion can be checked by an exact
retained local test. Take
\[
 f=x\bigl[y^2+xz(x-z)(x-2z)\bigr].
\]
It has only the smooth $A_3$ critical axis $x=y=0$. The
second normalization branch is the double plane branched
over four distinct lines. Blowing up their common point
and normalizing gives a double cover of the exceptional
$\mathbb P^1$ branched in four points, hence an elliptic
curve of square $-2$. Its discrepancy is $-1$ and the
order of $x$ on it is one, giving pole coefficient two
for the displayed total reduced hypersurface. It is excluded
here by the polar deformation, not by the coefficient alone.
For $H=z(x-z)(x-2z)$, the deformation $H\mapsto H-1$ has
the exact Groebner basis
\[
 3x-5z^4-14z,\qquad 25z^6+50z^3-2.
\]
It gives six simple solutions; adjoining $x$, $H_x$ or
$J$ to the critical ideal gives the unit ideal in each
case. For $H\mapsto H-\varepsilon$ all six scale by
$\varepsilon^{1/3}$ and converge to the origin. The
identities and these exact rational calculations are
executed in \texttt{code/part4/smooth\_a3\_polar.py},
with JSON and text output retained.
\end{proof}
